% Source: https://artofproblemsolving.com/community/c3h2808273_2011_geometry_bee__preliminary_round
% Stable identity: aops-topic-2808273-q07
Let $x$ be the diagonal separating the 52,39 sides from the 33,56 sides. The opposite angles have supplementary cosines. By the cosine rule their cosines are $\dfrac{4225-x^2}{2\cdot52\cdot39}$ and $\dfrac{4225-x^2}{2\cdot33\cdot56}$. Their sum is 0, forcing $x=65$. Both angles are right, so this diagonal is a diameter of the circumcircle. Every other chord has length at most the diameter; hence the longest diagonal is $\boxed{65}$.
\begin{asy}
usepackage("amsmath");
size(7cm);
pair A=(0,0),C=(65,0),B=(2704/65,2028/65),D=(3136/65,-1848/65);draw(circle((32.5,0),32.5));draw(A--B--C--D--cycle);draw(A--C,blue);label("A",A,W);label("B",B,N);label("C",C,E);label("D",D,S);label("52",(A+B)/2,NW);label("39",(B+C)/2,NE);label("33",(C+D)/2,SE);label("56",(D+A)/2,SW);
\end{asy}
